\section{Modeling setup} \label{sec:ModelingSetup}
The board uses standard transmission lines with an impedance of $Z_0=50\Omega$ and a time delay of $T_D=600ps$. The output impedance is assumed to be $11\Omega$, thus a series resistor of $39\Omega$ is added to match the output of the buffers with the transmission lines.


\subsection{Draw models} \label{subsec:models}
The model that we need to draw has to include all the parasitics of the pins of the IC's. Each pin has both a series resistance and inductance that need to be accounted for, but also two parallel capacitors to $GND$, one on the pin side and one from the die side. Lastly we need to have the internal capacitances taken account of. These consist of the capacitance to the die $VCC$ and $GND$. On the other end of the transmission line we assume that there is a capacitance on the input pins of the receiving buffer. The internal part of the driver is modeled as to resistors controlled by the input clock. For the passive line it is either on or off. 

The passive line is modeled in figure \ref{fig:monLine}, and the first Active line is modeled in figure \ref{fig:ActLine}.

\begin{figure}[h]
    \centering
    \includegraphics[width=\linewidth]{MonitorLine.png}
    \caption{The monitored line is modeled with all the parameters needed to simulate it. Without values}
    \label{fig:monLine}
\end{figure}

\begin{figure}[h]
    \centering
    \includegraphics[width=\linewidth]{ActiveLineModel.png}
    \caption{The Active line L1 is modeled with all the parameters needed to simulate it. Without values}
    \label{fig:ActLine}
\end{figure}

\subsection{Estimate parameters} \label{subsec:parameters}
All the parameters that need to be estimated are either found in the datasheet or the ibis file for the SN74ALVC244, since all other values are given. Or are in the schematic.

The values from the schematic that are used is the decoupling capacitor which is set to $C_{dec}=100nF$ here we don't take the capacitor's parasitics into account. The die capacitance is common for both packages since the die is the same, this capacitance is $C_{comp}=6.04pF$, here it is important to not that this is to both lines next to the given line, so the values of $Cup$ and $Cdn$ are half of $C_{comp}$, so $Cup=Cdn=C_{comp}/2=3.02pF$. The off resistance of the buffer is assumed to be $R_{off}=1M\Omega$.

The receiver capacitance for all the active lines are defined as the $C_{comp}=4.21pF$ value that is constant for both the SOIC and the QFN packages, plus the $C_{pin}$ value for each pin on the receiver.

 
\begin{table}[htbp]
    \centering
    \begin{tabular}{|c|c|c|c|c|c|c|c|}
        \hline
        & & \multicolumn{3}{c|}{\textbf{QFN (RGY)}} & \multicolumn{3}{c|}{\textbf{SOIC (DW)}} \\
        \hline
        \textbf{Pin} & \textbf{Signal} & $R$ (m$\Omega$) & $L$ (nH) & $C$ (pF) & $R$ (m$\Omega$) & $L$ (nH) & $C$ (pF) \\
        \hline
        20 & VCC & 91.31 & 1.3037 & 0.12549 & 42 & 6.308 & 0.864 \\
        10 & GND & 45.79 & 1.00272 & 0.15399 & 43 & 6.349 & 0.866 \\
        3 & 2Y4 & 101.96 & 1.47947 & 0.14411 & 34 & 4.714 & 0.732 \\
        5 & 2Y3 & 93.15 & 1.34070 & 0.12663 & 36 & 4.130 & 0.533 \\
        18 & 1Y1 & 97.27 & 1.39536 & 0.13843 & 34 & 4.714 & 0.731 \\
        16 & 1Y2 & 93.36 & 1.33430 & 0.12780 & 36 & 4.130 & 0.533 \\
        14 & 1Y3 & 100.13 & 1.44133 & 0.14031 & 35 & 4.414 & 0.596 \\
        12 & 1Y4 & 118.19 & 1.73153 & 0.15918 & 45 & 5.561 & 0.868 \\
        9 & 2Y1 & 110.53 & 1.61643 & 0.15073 & 40 & 5.518 & 0.865 \\
        7 & 2Y2 & 95.98 & 1.38423 & 0.13711 & 34 & 4.284 & 0.583 \\
        \hline
    \end{tabular}
    \caption{Package parasitics per pin from the IBIS file (components ALVC244\_RGY and ALVC244\_DW). Each $C_{pin}$ is split as $C_{pin}/2$ at the die side and $C_{pin}/2$ at the board side; for the $GND$ pin only the die-side half is used.}
    \label{tab:sim_package}
\end{table}




\subsection{Noise margins} \label{subsec:noise-margins}

\[ V_{CC,\mathrm{nom}}=3.3V \qquad |I_O| \leq 12mA \]

\begin{table}[h]
    \centering
    \begin{tabular}{ccc}
         \hline
         Parameters & Values & Datasheet conditions\\
         \hline
         $V_{OH,\min}$ & $2.4V$ & $V_{CC}=3.0V$, $I_{OH}=-12mA$ \\
         $V_{OL,\max}$ & $0.4V$ & $V_{CC}=2.7V$, $I_{OL}=+12mA$ \\
         $V_{IH,\min}$ & $2.0V$ & $V_{CC}=2.7$--$3.6V$ \\
         $V_{IL,\max}$ & $0.8V$ & $V_{CC}=2.7$--$3.6V$ \\
         \hline
    \end{tabular}
    \caption{SN74ALVC244 logic levels used in the noise-margin calculation.}
    \label{tab:logic-levels}
\end{table}

The higher and lower noise margins can be calculated from the values of \ref{tab:logic-levels}. The two levels are calculated in eq \ref{eq:noise-margin-high} and \ref{eq:noise-margin-low}.

\begin{equation}
NM_H=V_{OH,\min}-V_{IH,\min}=2.4V-2.0V=0.4V.
\label{eq:noise-margin-high}
\end{equation}

\begin{equation}
NM_L=V_{IL,\max}-V_{OL,\max}=0.8V-0.4V=0.4V.
\label{eq:noise-margin-low}
\end{equation}

We compare the driver's output levels with the receiver's input thresholds. The voltage levels were obtained from the datasheet, and we used formulas to calculate the difference between them—the margins. The results yield what are known as noise margins, which in both cases amount to 0.4 volts. These margins indicate the extent to which the voltage can drop or rise due to noise while still ensuring that the receiver can reliably distinguish between zeros and ones. If voltage fluctuations exceed these margins—that is, if the voltage for a 'one' drops by more than 0.4 volts, or the voltage for a 'zero' rises by more than $0.4V$ (relative to the worst-case output levels)—then correct recognition of the logic level can no longer be guaranteed.


